\documentclass[12pt]{article}

\usepackage[T1]{fontenc}
\usepackage{lmodern}
\usepackage[margin=1in]{geometry}
\usepackage{newpxtext,newpxmath}
\usepackage{amsmath}
\usepackage{parskip}

\linespread{1.03}

\title{\Large Solutions to Practice Final Exam \\[0.3em]}
\author{}
\date{}

\begin{document}

\maketitle
\vspace{-5em}

\begin{table}[h!]
    \centering
    \renewcommand{\arraystretch}{1.2}
    \begin{tabular}{cc @{\hspace{1cm}} cc @{\hspace{1cm}} cc}
        \hline
        \textbf{Problem} & \textbf{Answer} & \textbf{Problem} & \textbf{Answer} & \textbf{Problem} & \textbf{Answer} \\
        \hline
        1                & B               & 9                & B               & 17               & D               \\
        2                & D               & 10               & C               & 18               & A               \\
        3                & D               & 11               & B               & 19               & C               \\
        4                & A               & 12               & B               & 20               & B               \\
        5                & C               & 13               & C               & 21               & B               \\
        6                & D               & 14               & D               & 22               & B               \\
        7                & D               & 15               & C               & 23               & B               \\
        8                & D               & 16               & A               &                  &                 \\
        \hline
    \end{tabular}
\end{table}

\section*{Problem 24}
\emph{In general, a clear and complete argument received full credit. An argument that missed a key point received half credit (-2).}

\begin{itemize}
    \item[a)] It is efficient for the professor to assign the new textbook. The additional surplus associated with the new textbook is \$2000 (1000 students, \$2 each). The surplus associated with the Classic Textbook is \$1000, accruing to the Author.

    \item[]
          (The professor and Morris do not benefit one way or the other, so they are not of concern.)

    \item[b)]
          The Author is highly incentivized to influence (i.e., lobby) the professor, for example by taking him out to dinner and persuading the professor to keep the Classic Textbook.

          Each student benefits only \$2. No individual student has a strong incentive to influence the professor, and it would be difficult for students to organize in a way that could effectively compete with the Author.

          As a result, the Author is likely to succeed in swaying the professor to keep the Classic Textbook.

          (Note: Discussing switching costs or other potential barriers to entry was a good idea, but by itself was not sufficient to earn \emph{full} credit. To receive full credit, it was necessary to identify the potential for lobbying that is central to this problem.)

    \item[c)]
          If students could organize and act collectively, they would have an incentive to lobby or otherwise influence the professor.

          In general, organizing is costly; however, if the cost of organizing were sufficiently low, students would find it beneficial to do so.

          (The professor does not care one way or the other. Morris also does not care one way or the other.)
\end{itemize}

\newpage
\section*{Problem 25}
The market demand and cost information was:
\[
    Q = 10 - P, \quad MR = 10 - 2Q, \quad MC = 6, \quad AC = 6 + \frac{2}{Q}.
\]

\subsubsection*{(a) Efficient Quantity}

Efficiency requires setting price equal to marginal cost:
\[
    P = MC = 6.
\]
Using the demand curve,
\[
    Q = 10 - 6 = 4.
\]

\[
    \boxed{Q^{\text{eff}} = 4}
\]

\subsubsection*{(b) Profit-Maximizing Price}

Profit maximization requires \(MR = MC\):
\[
    10 - 2Q = 6 \quad \Rightarrow \quad Q^* = 2.
\]
Substituting into the demand curve,
\[
    P^* = 10 - 2 = 8.
\]

\[
    \boxed{P^* = 8}
\]

\subsubsection*{(c) Deadweight Loss}

Deadweight loss is the triangle between \(Q^*\) and \(Q^{\text{eff}}\):
\[
    \text{DWL} = \frac{1}{2}(Q^{\text{eff}} - Q^*)(P^* - MC)
    = \frac{1}{2}(4 - 2)(8 - 6) = 2.
\]

\[
    \boxed{\text{DWL} = 2}
\]

\subsubsection*{(d) Profit}

Profit equals \((P^* - AC)\times Q^*\). First compute average cost:
\[
    AC = 6 + \frac{2}{2} = 7.
\]
Thus,
\[
    \pi = (8 - 7)\times 2 = 2.
\]

\[
    \boxed{\pi = 2}
\]


\end{document}